\section{Related Work}
\label{sec:related}

\subsection{Pedestrian Attribute Recognition (PAR)}
Pedestrian Attribute Recognition has evolved from early handcrafted feature pipelines to deep end-to-end multi-label learning systems. Early landmark approaches, such as HydraPlus-Net~\cite{liu2017hydraplus}, incorporated multi-directional attention modules to aggregate low-level and high-level feature maps. Tang \etal~\cite{tang2019improving} introduced weakly-supervised attention mechanisms to locate attribute regions without manual bounding box annotations. Graph-based formulations, such as Visual-Semantic Graph Reasoning~\cite{li2019visual}, model semantic correlations between attributes via graph convolutional networks.

A major milestone in benchmarking PAR under realistic surveillance conditions was established by the UPAR dataset and challenges~\cite{specker2023upar,cormier2023upar,cormier2024upar}. Specker \etal~\cite{specker2023upar} harmonized diverse public benchmarks---including Market1501~\cite{zheng2015scalable}, PA-100K~\cite{liu2017hydraplus}, and PETA~\cite{deng2014pedestrian}---into a unified 40-attribute taxonomy, demonstrating that cross-domain generalization and consistent annotation protocols are essential for real-world surveillance deployment.

\subsection{Modern Vision Backbones}
While earlier PAR models predominantly relied on ResNet-50~\cite{he2016deep} backbones, recent architectural advancements have revitalized convolutional neural networks. ConvNeXt~\cite{liu2022convnet} modernizes standard ConvNets by incorporating design principles from vision transformers, including depthwise separable $7 \times 7$ convolutions, inverted bottleneck ratios, fewer activation functions (GELU), and LayerNorm. ConvNeXt preserves the translational equivariance and spatial inductive bias of convolutions while matching the receptive field and performance of Vision Transformers, making it an ideal feature backbone for pedestrian analysis.

\subsection{Spatial Strip Pooling and Part Models}
Due to the vertical elongation and canonical upright orientation of pedestrians in surveillance cameras, horizontal strip pooling has proven effective in person re-identification. Sun \etal~\cite{sun2018beyond} introduced the Part-based Convolutional Baseline (PCB), slicing feature maps into uniform horizontal stripes to capture spatial structure. However, standard PCB feeds all horizontal stripes into uniform classifiers without semantic distinction. In contrast, our ASR-HC framework couples horizontal anatomical partitioning directly with domain-specific attribute semantics: strip tokens are refined through self-attention across anatomical positions and explicitly routed to body-specific attribute sub-spaces (e.g., hats/glasses to the upper strip, trousers/skirts to the lower strip).

\subsection{Class Imbalance and Logit Calibration}
Class imbalance is ubiquitous in real-world multi-label vision tasks. Menon \etal~\cite{menon2021longtail} established that Fisher-consistent optimization under class-imbalanced data can be achieved via logit adjustment using training class priors. In multi-label surveillance PAR, this issue is exacerbated by the diverse frequencies of 40 attributes: ubiquitous demographic attributes dominate optimization, while rare clothing colors suffer high false-negative rates. In this work, we propose a Source-Prior Harmonic Calibrator that leverages empirical training priors $\pi_j$ combined with competitive margin suppression across mutually exclusive color classes, resolving color hallucinations while optimizing the Harmonic Mean.
